
%%%%%%%%%%%%%%%%%%%%%%%%%%%%%%%%%%%%%%%%%%%%%%%%%%%%%%%%%%%%%%%%%%%
% Physik IV - Blatt 11 Lernfassung / Abgabe / Tafelvortrag
% Standalone university LaTeX template -- one-file version
%%%%%%%%%%%%%%%%%%%%%%%%%%%%%%%%%%%%%%%%%%%%%%%%%%%%%%%%%%%%%%%%%%%

\documentclass[11pt,a4paper]{scrartcl}

%%%%%%%%%%%%%%%%%%%%%%%%%%%%%%%%%%%%%%%%%%%%%%%%%%%%%%%%%%%%%%%%%%%
% Embedded file: includes/university_metadata.tex
%%%%%%%%%%%%%%%%%%%%%%%%%%%%%%%%%%%%%%%%%%%%%%%%%%%%%%%%%%%%%%%%%%%
\newcommand{\CourseTitle}{COURSE NAME}
\newcommand{\CourseShortTitle}{COURSE SHORT NAME}
\newcommand{\DocumentKind}{DOCUMENT KIND}
\newcommand{\DocumentTitle}{DOCUMENT TITLE}
\newcommand{\DocumentSubtitle}{}
\newcommand{\AuthorName}{YOUR NAME}
\newcommand{\InstitutionName}{INSTITUTION NAME}
\newcommand{\SubmissionDate}{\today}

\newcommand{\makeuniversitytitle}{%
    \begin{titlepage}
        \centering
        \vspace*{2cm}
        {\Large \CourseTitle\par}
        \vspace{1.2cm}
        {\huge\bfseries \DocumentTitle\par}
        \ifx\DocumentSubtitle\empty\else
            \vspace{0.4cm}
            {\Large \DocumentSubtitle\par}
        \fi
        \vspace{1cm}
        {\Large \DocumentKind\par}
        \vfill
        {\large \AuthorName\par}
        \vspace{0.2cm}
        {\large \InstitutionName\par}
        \vspace{0.8cm}
        {\large \SubmissionDate\par}
    \end{titlepage}
    \pagenumbering{arabic}
}

\newcommand{\makechaptertitle}{%
    \begin{center}
        {\Large \CourseTitle\par}
        \vspace{0.5cm}
        {\huge\bfseries \DocumentTitle\par}
        \ifx\DocumentSubtitle\empty\else
            \vspace{0.25cm}
            {\Large \DocumentSubtitle\par}
        \fi
        \vspace{0.5cm}
        {\large \AuthorName \hfill \SubmissionDate\par}
    \end{center}
    \vspace{0.5cm}
}

\newcommand{\makephysiktitle}{\makeuniversitytitle}

%%%%%%%%%%%%%%%%%%%%%%%%%%%%%%%%%%%%%%%%%%%%%%%%%%%%%%%%%%%%%%%%%%%
% Document metadata
%%%%%%%%%%%%%%%%%%%%%%%%%%%%%%%%%%%%%%%%%%%%%%%%%%%%%%%%%%%%%%%%%%%
\renewcommand{\CourseTitle}{Physik IV: Kerne und Teilchen}
\renewcommand{\CourseShortTitle}{Physik IV}
\renewcommand{\DocumentKind}{Hausaufgabe, Tafelvortrag und Lernfassung}
\renewcommand{\DocumentTitle}{Aufgabenblatt 11}
\renewcommand{\DocumentSubtitle}{Z-Resonanz und solare Neutrinos}
\renewcommand{\AuthorName}{Tobias Laser}
\renewcommand{\InstitutionName}{Universit\"at Innsbruck}
\renewcommand{\SubmissionDate}{6. Juni 2026}

%%%%%%%%%%%%%%%%%%%%%%%%%%%%%%%%%%%%%%%%%%%%%%%%%%%%%%%%%%%%%%%%%%%
% Embedded file: includes/university_packages.tex
%%%%%%%%%%%%%%%%%%%%%%%%%%%%%%%%%%%%%%%%%%%%%%%%%%%%%%%%%%%%%%%%%%%
\usepackage[margin=2.5cm]{geometry}
\usepackage[onehalfspacing]{setspace}
\usepackage{scrlayer-scrpage}

\usepackage[T1]{fontenc}
\usepackage[utf8]{inputenc}
\usepackage[ngerman]{babel}
\usepackage{lmodern}
\usepackage{microtype}
\usepackage{csquotes}

\usepackage{amsmath,amsthm,amssymb,mathtools}
\usepackage{bm}
\usepackage{icomma}
\usepackage{cancel}

\usepackage{graphicx}
\usepackage{tikz}
\usetikzlibrary{arrows.meta,calc,decorations.pathmorphing,positioning,patterns,angles,quotes}
\usepackage{pgfplots}
\pgfplotsset{compat=1.18}

\usepackage[locale=DE,space-before-unit=true,per-mode=symbol,separate-uncertainty=true]{siunitx}
\usepackage{booktabs,multirow,array,tabularx,longtable}
\usepackage{caption,subcaption}
\usepackage{placeins}
\usepackage{enumitem}
\usepackage{xparse}

\usepackage[most]{tcolorbox}
\tcbuselibrary{breakable,skins,theorems}

\usepackage[breaklinks=true,colorlinks=true,linkcolor=blue,urlcolor=blue,citecolor=blue]{hyperref}
\usepackage[nameinlink,noabbrev]{cleveref}

\clearpairofpagestyles
\ihead{\CourseShortTitle}
\ohead{\DocumentKind}
\cfoot{\pagemark}
\pagestyle{scrheadings}

\setlist[enumerate]{itemsep=0.4em,topsep=0.4em}
\setlist[itemize]{itemsep=0.25em,topsep=0.35em}
\captionsetup{font=small,labelfont=bf}

%%%%%%%%%%%%%%%%%%%%%%%%%%%%%%%%%%%%%%%%%%%%%%%%%%%%%%%%%%%%%%%%%%%
% Embedded file: includes/university_macros.tex
%%%%%%%%%%%%%%%%%%%%%%%%%%%%%%%%%%%%%%%%%%%%%%%%%%%%%%%%%%%%%%%%%%%
\newcommand{\R}{\mathbb{R}}
\newcommand{\N}{\mathbb{N}}
\newcommand{\Z}{\mathbb{Z}}
\newcommand{\Q}{\mathbb{Q}}
\newcommand{\C}{\mathbb{C}}

\newcommand{\set}[1]{\left\{ #1 \right\}}
\newcommand{\abs}[1]{\left| #1 \right|}
\newcommand{\norm}[1]{\left\lVert #1 \right\rVert}
\newcommand{\avg}[1]{\left\langle #1 \right\rangle}
\newcommand{\paren}[1]{\left( #1 \right)}
\newcommand{\bracket}[1]{\left[ #1 \right]}
\newcommand{\ol}[1]{\overline{#1}}

\newcommand{\dd}{\,\mathrm{d}}
\newcommand{\dv}[2]{\frac{\dd #1}{\dd #2}}
\newcommand{\pdv}[2]{\frac{\partial #1}{\partial #2}}
\newcommand{\pddv}[2]{\frac{\partial^2 #1}{\partial #2^2}}
\newcommand{\evalat}[2]{\left. #1 \right|_{#2}}

\newcommand{\vect}[1]{\bm{#1}}
\newcommand{\unitvect}[1]{\hat{\bm{#1}}}
\newcommand{\grad}{\vect{\nabla}}
\newcommand{\divergence}{\grad\!\cdot}
\newcommand{\curl}{\grad\!\times}

\newcommand{\half}{\frac{1}{2}}
\newcommand{\third}{\frac{1}{3}}
\newcommand{\twothirds}{\frac{2}{3}}
\newcommand{\hbarc}{\hbar c}
\newcommand{\alphaf}{\alpha}
\newcommand{\eV}{\si{eV}}
\newcommand{\MeV}{\si{MeV}}
\newcommand{\GeV}{\si{GeV}}
\newcommand{\TeV}{\si{TeV}}

\newcommand{\nuclide}[3]{^{#1}_{#2}\mathrm{#3}}
\newcommand{\isotope}[2]{^{#1}\mathrm{#2}}
\newcommand{\anti}[1]{\overline{#1}}
\newcommand{\pbar}{\anti{p}}
\newcommand{\nbar}{\anti{n}}
\newcommand{\ee}{e^-}
\newcommand{\ep}{e^+}
\newcommand{\mup}{\mu^+}
\newcommand{\mum}{\mu^-}
\newcommand{\nue}{\nu_e}
\newcommand{\numu}{\nu_\mu}
\newcommand{\nutau}{\nu_\tau}
\newcommand{\pionp}{\pi^+}
\newcommand{\pionm}{\pi^-}
\newcommand{\pionz}{\pi^0}
\newcommand{\kaonp}{K^+}
\newcommand{\kaonm}{K^-}
\newcommand{\kaonz}{K^0}

\newcommand{\diffxs}{\frac{\dd\sigma}{\dd\Omega}}
\newcommand{\totalxs}{\sigma_{\mathrm{total}}}
\newcommand{\lum}{\mathcal{L}}
\newcommand{\smand}{s}
\newcommand{\tmand}{t}
\newcommand{\umand}{u}
\newcommand{\Ecm}{E_{\mathrm{CM}}}

\newcommand{\ie}{d.\,h.}
\newcommand{\eg}{z.\,B.}
\newcommand{\wrt}{bez\"uglich}

\DeclareSIUnit{\barn}{b}
\DeclareSIUnit{\millibarn}{mb}
\DeclareSIUnit{\microbarn}{\micro b}
\DeclareSIUnit{\nanobarn}{nb}
\DeclareSIUnit{\fermi}{fm}
\DeclareSIUnit{\year}{a}

%%%%%%%%%%%%%%%%%%%%%%%%%%%%%%%%%%%%%%%%%%%%%%%%%%%%%%%%%%%%%%%%%%%
% Embedded file: includes/university_boxes.tex
%%%%%%%%%%%%%%%%%%%%%%%%%%%%%%%%%%%%%%%%%%%%%%%%%%%%%%%%%%%%%%%%%%%
\tcbset{
    unibox/.style={
        enhanced,
        breakable,
        sharp corners,
        boxrule=0.6pt,
        left=1.2mm,
        right=1.2mm,
        top=1mm,
        bottom=1mm,
        before skip=0.8em,
        after skip=0.8em,
        fonttitle=\bfseries
    }
}

\newtcolorbox{calculationbox}[1][]{unibox,title=Rechnung,colback=black!2,colframe=black!55,#1}
\newtcolorbox{sidecalcbox}[1][]{unibox,title=Nebenrechnung,colback=black!1,colframe=black!40,#1}
\newtcolorbox{solutionbox}[1][]{unibox,title=L\"osung,colback=blue!2,colframe=blue!45!black,#1}
\newtcolorbox{answerbox}[1][]{unibox,title=Antwort,colback=green!3,colframe=green!45!black,#1}
\newtcolorbox{notebox}[1][]{unibox,title=Notiz,colback=yellow!6,colframe=yellow!45!black,#1}
\newtcolorbox{definitionbox}[1][]{unibox,title=Definition,colback=cyan!3,colframe=cyan!45!black,#1}
\newtcolorbox{warningbox}[1][]{unibox,title=Achtung,colback=red!3,colframe=red!55!black,#1}
\newtcolorbox{summarybox}[1][]{unibox,title=Zusammenfassung,colback=violet!3,colframe=violet!50!black,#1}
\newtcolorbox{formulabox}[1][]{unibox,title=Formel,colback=orange!4,colframe=orange!65!black,#1}
\newtcolorbox{taskbox}[1][]{unibox,title=Aufgabe,colback=white,colframe=black!75,#1}
\newtcolorbox{talkbox}[1][]{unibox,title=Tafelvortrag,colback=teal!3,colframe=teal!50!black,#1}
\newtcolorbox{learningbox}[1][]{unibox,title=Lernidee,colback=purple!3,colframe=purple!50!black,#1}

%%%%%%%%%%%%%%%%%%%%%%%%%%%%%%%%%%%%%%%%%%%%%%%%%%%%%%%%%%%%%%%%%%%
% Embedded file: includes/university_theorems.tex
%%%%%%%%%%%%%%%%%%%%%%%%%%%%%%%%%%%%%%%%%%%%%%%%%%%%%%%%%%%%%%%%%%%
\theoremstyle{definition}
\newtheorem{definition}{Definition}[section]
\newtheorem{example}{Beispiel}[section]

\theoremstyle{plain}
\newtheorem{satz}{Satz}[section]
\newtheorem{lemma}{Lemma}[section]

\theoremstyle{remark}
\newtheorem*{remark}{Bemerkung}

%%%%%%%%%%%%%%%%%%%%%%%%%%%%%%%%%%%%%%%%%%%%%%%%%%%%%%%%%%%%%%%%%%%
\begin{document}

\makeuniversitytitle
\tableofcontents
\newpage

%%%%%%%%%%%%%%%%%%%%%%%%%%%%%%%%%%%%%%%%%%%%%%%%%%%%%%%%%%%%%%%%%%%
\section*{Einordnung des Blattes}
\addcontentsline{toc}{section}{Einordnung des Blattes}

\begin{summarybox}[title=Ziel dieses Blattes]
Dieses Blatt verbindet zwei typische Rechnungen der Teilchen- und Astroteilchenphysik:
\begin{itemize}
    \item Bei der \textbf{$Z$-Resonanz} sieht man, wie ein instabiles Austauschteilchen als Resonanz im Wirkungsquerschnitt erscheint.
    \item Bei den \textbf{solaren Neutrinos} sieht man, warum ein riesiger Teilchenfluss wegen winziger Wirkungsquerschnitte trotzdem kaum Wechselwirkungen im Körper erzeugt.
\end{itemize}
Der wichtigste konzeptionelle Punkt ist: Ein Wirkungsquerschnitt ist eine effektive Trefferfläche. Große Flüsse allein bedeuten nicht automatisch große Reaktionsraten; entscheidend ist immer das Produkt $\Phi\sigma N$.
\end{summarybox}

\begin{notebox}[title=Benutzte Konstanten und Konventionen]
\[
    \hbar c = \SI{197}{MeV.fm}=\SI{0.197}{GeV.fm},
    \qquad
    \alpha=\frac{1}{137.036},
    \qquad
    \SI{1}{barn}=\SI{e-28}{m^2}.
\]
\[
    \SI{1}{fm^2}=\SI{e-30}{m^2}=10^{-2}\,\si{barn}=10^7\,\si{nb},
    \qquad
    \SI{1}{eV}=\SI{1.602e-19}{J}.
\]
\end{notebox}

\FloatBarrier
\newpage

%%%%%%%%%%%%%%%%%%%%%%%%%%%%%%%%%%%%%%%%%%%%%%%%%%%%%%%%%%%%%%%%%%%
\section{Aufgabe 1: Z-Resonanz}

\begin{taskbox}
Am Large-Electron-Positron Collider (LEP) wurden Elektronen und Positronen so beschleunigt, dass im Schwerpunktsystem die Ruheenergie des $Z$-Bosons zur Verfügung stand:
\[
    E_{\mathrm{CM}}=m_Zc^2=\SI{91.2}{GeV}.
\]
Bei dieser Energie ist der Wirkungsquerschnitt für die Produktion des $Z$-Bosons maximal. Der differentielle Wirkungsquerschnitt sei
\[
    \diffxs(E_{\mathrm{CM}},\theta)
    =(\hbar c)^2\frac{\alpha^2}{s}\frac{G(s)}{4}\paren{1+\cos^2\theta},
\]
mit
\[
    G(s)=1+\frac{s^2g_eg_\mu}{\paren{s-m_Z^2c^4}^2+m_Z^2c^4\Gamma_Z^2},
    \qquad
    s=E_{\mathrm{CM}}^2,
\]
\[
    \Gamma_Z=\SI{2.5}{GeV},
    \qquad
    g_e=g_\mu=0.355.
\]
Der erste Term in $G(s)$ beschreibt den Beitrag über ein virtuelles Photon, der zweite Term den Beitrag über das $Z$-Boson.
\begin{enumerate}[label=(\alph*)]
    \item Um welchen Faktor ist der zum $Z$-Boson gehörige Term am Maximum der Resonanz größer als der Beitrag des virtuellen Photons?
    \item Skizzieren Sie die Winkelabhängigkeit des differentiellen Wirkungsquerschnitts. Vergleichen Sie mit der entsprechenden Abhängigkeit des Rutherfordstreuquerschnitts.
    \item Berechnen Sie den totalen Wirkungsquerschnitt $\sigma_{\mathrm{total}}$, indem Sie den differentiellen Wirkungsquerschnitt über den gesamten Raumwinkel integrieren.
    \item Wie groß ist somit der totale Wirkungsquerschnitt bei $E_{\mathrm{CM}}=m_Zc^2$ in Nanobarn? Welchen Radius hätte ein Kreis mit gleicher Fläche?
    \item Plotten Sie $\sigma_{\mathrm{total}}$ als Funktion von $E_{\mathrm{CM}}$ im Bereich \SIrange{50}{150}{GeV} linear und im Bereich \SIrange{10}{500}{GeV} doppeltlogarithmisch.
\end{enumerate}
\end{taskbox}

\subsection{Lernteil: Was bedeutet eine Resonanz im Wirkungsquerschnitt?}

\begin{learningbox}
Eine Resonanz entsteht, wenn die verfügbare Schwerpunktenergie genau zur Ruheenergie eines instabilen Zwischenzustands passt. Hier ist das der Fall bei
\[
    E_{\mathrm{CM}}\approx m_Zc^2.
\]
Dann kann in der Reaktion $e^+e^-\to Z\to \mu^+\mu^-$ das $Z$-Boson resonant erzeugt werden. Mathematisch sieht man das am Nenner des $Z$-Terms in $G(s)$: Bei $s=m_Z^2c^4$ verschwindet der Term $(s-m_Z^2c^4)^2$. Übrig bleibt im Nenner nur die Breite $\Gamma_Z$, die den Peak endlich macht.
\end{learningbox}

\subsection{Abgabe: Ausführliche Lösung}

\begin{enumerate}[label=(\alph*)]
    \item \textbf{Resonanzfaktor.}

    \begin{calculationbox}
        Am Maximum der Resonanz gilt
        \[
            E_{\mathrm{CM}}=m_Zc^2
            \quad\Longrightarrow\quad
            s=E_{\mathrm{CM}}^2=m_Z^2c^4.
        \]
        Damit verschwindet im Nenner des $Z$-Terms der Ausdruck
        \[
            \paren{s-m_Z^2c^4}^2=0.
        \]
        Der Photonbeitrag in $G(s)$ ist
        \[
            G_\gamma=1.
        \]
        Der $Z$-Beitrag am Resonanzmaximum ist
        \[
            G_Z
            =\frac{s^2g_eg_\mu}{m_Z^2c^4\Gamma_Z^2}.
        \]
        Da am Maximum $s=m_Z^2c^4$ gilt, folgt
        \[
            G_Z=\frac{s g_eg_\mu}{\Gamma_Z^2}.
        \]
        Mit
        \[
            s=(\SI{91.2}{GeV})^2
        \]
        erhält man
        \[
            G_Z=\frac{(91.2)^2\cdot0.355\cdot0.355}{(2.5)^2}=167.7.
        \]
    \end{calculationbox}

    \begin{answerbox}
        Am Resonanzmaximum ist der $Z$-Beitrag ungefähr
        \[
            \boxed{\frac{G_Z}{G_\gamma}\approx 1.68\cdot10^2}
        \]
        mal größer als der Beitrag des virtuellen Photons.
    \end{answerbox}

    \item \textbf{Winkelabhängigkeit.}

    \begin{calculationbox}
        Für festes $E_{\mathrm{CM}}$ sind alle Faktoren bis auf den Winkel konstant. Daher ist
        \[
            \diffxs\propto 1+\cos^2\theta.
        \]
        Die Funktion ist symmetrisch unter $\theta\mapsto \pi-\theta$. Sie ist maximal bei
        \[
            \theta=0,\pi,
            \qquad
            1+\cos^2\theta=2,
        \]
        und minimal bei
        \[
            \theta=\frac{\pi}{2},
            \qquad
            1+\cos^2\theta=1.
        \]
    \end{calculationbox}

    \begin{calculationbox}[title=Skizze]
    \begin{center}
    \begin{tikzpicture}
    \begin{axis}[
        width=0.84\textwidth,
        height=5.8cm,
        xlabel={$\theta$ in Grad},
        ylabel={relative Stärke},
        xmin=0,xmax=180,
        ymin=0,ymax=2.2,
        xtick={0,45,90,135,180},
        ytick={0,1,2},
        grid=both,
        legend pos=south east
    ]
        \addplot[thick,domain=0:180,samples=200]{1+cos(x)^2};
        \addlegendentry{$1+\cos^2\theta$}
    \end{axis}
    \end{tikzpicture}
    \end{center}
    \end{calculationbox}

    \begin{calculationbox}[title=Vergleich mit Rutherfordstreuung]
        Für die Rutherfordstreuung gilt qualitativ
        \[
            \diffxs_R\propto \frac{1}{\sin^4(\theta/2)}.
        \]
        Für $\theta\to0$ divergiert dieser Ausdruck. Rutherfordstreuung ist daher stark vorwärtsgerichtet. Die Verteilung $1+\cos^2\theta$ ist dagegen überall endlich und ändert sich nur um einen Faktor $2$ zwischen Minimum und Maximum.
    \end{calculationbox}

    \begin{answerbox}
        \[
            \boxed{\diffxs\propto 1+\cos^2\theta}
        \]
        Die Winkelverteilung ist glatt, endlich und vorwärts-rückwärts-symmetrisch. Rutherfordstreuung ist dagegen wegen $1/\sin^4(\theta/2)$ stark vorwärtsdominant und divergiert für kleine Winkel.
    \end{answerbox}

    \item \textbf{Integration zum totalen Wirkungsquerschnitt.}

    \begin{calculationbox}
        Der totale Wirkungsquerschnitt ist
        \[
            \sigma_{\mathrm{total}}=\int \diffxs\dd\Omega.
        \]
        Mit
        \[
            \dd\Omega=\sin\theta\dd\theta\dd\varphi,
            \qquad
            \theta\in[0,\pi],
            \quad
            \varphi\in[0,2\pi],
        \]
        folgt
        \[
            \sigma_{\mathrm{total}}
            =(\hbar c)^2\frac{\alpha^2}{s}\frac{G(s)}{4}
            \int_0^{2\pi}\int_0^\pi (1+\cos^2\theta)\sin\theta\dd\theta\dd\varphi.
        \]
        Das $\varphi$-Integral ist
        \[
            \int_0^{2\pi}\dd\varphi=2\pi.
        \]
        Für das $\theta$-Integral setzen wir $u=\cos\theta$, also $\dd u=-\sin\theta\dd\theta$:
        \[
        \begin{aligned}
            \int_0^\pi (1+\cos^2\theta)\sin\theta\dd\theta
            &=\int_1^{-1}(1+u^2)(-\dd u)\\
            &=\int_{-1}^{1}(1+u^2)\dd u\\
            &=\bracket{u+\frac{u^3}{3}}_{-1}^{1}=\frac{8}{3}.
        \end{aligned}
        \]
        Insgesamt ist also
        \[
            \int(1+\cos^2\theta)\dd\Omega
            =2\pi\cdot\frac{8}{3}=\frac{16\pi}{3}.
        \]
        Damit folgt
        \[
            \sigma_{\mathrm{total}}
            =(\hbar c)^2\frac{\alpha^2}{s}\frac{G(s)}{4}\cdot\frac{16\pi}{3}
            =(\hbar c)^2\frac{\alpha^2}{s}G(s)\frac{4\pi}{3}.
        \]
    \end{calculationbox}

    \begin{answerbox}
        \[
            \boxed{\sigma_{\mathrm{total}}(s)=(\hbar c)^2\frac{\alpha^2}{s}G(s)\frac{4\pi}{3}}
        \]
    \end{answerbox}

    \item \textbf{Numerischer Wert am Resonanzmaximum und Kreisradius.}

    \begin{calculationbox}
        Am Maximum verwenden wir
        \[
            G(s)=1+167.7=168.7,
            \qquad
            s=(\SI{91.2}{GeV})^2=\SI{8317.44}{GeV^2}.
        \]
        Mit $\hbar c=\SI{0.197}{GeV.fm}$ und $\alpha=1/137.036$ ergibt sich
        \[
        \begin{aligned}
            \sigma_{\mathrm{total}}
            &=(0.197)^2\,\si{fm^2}\frac{(1/137.036)^2}{8317.44}\cdot168.7\cdot\frac{4\pi}{3}\\
            &=1.756\cdot10^{-7}\,\si{fm^2}.
        \end{aligned}
        \]
        Wegen
        \[
            \SI{1}{fm^2}=10^7\,\si{nb}
        \]
        folgt
        \[
            \sigma_{\mathrm{total}}=1.756\cdot10^{-7}\cdot10^7\,\si{nb}=\SI{1.756}{nb}.
        \]
        Ein Kreis gleicher Fläche erfüllt
        \[
            \sigma_{\mathrm{total}}=\pi r^2.
        \]
        Also
        \[
            r=\sqrt{\frac{\sigma_{\mathrm{total}}}{\pi}}
            =\sqrt{\frac{1.756\cdot10^{-7}\,\si{fm^2}}{\pi}}
            =2.36\cdot10^{-4}\,\si{fm}.
        \]
        In Metern ist
        \[
            r=2.36\cdot10^{-19}\,\si{m}.
        \]
    \end{calculationbox}

    \begin{answerbox}
        Am Resonanzmaximum ergibt sich
        \[
            \boxed{\sigma_{\mathrm{total}}=\SI{1.76}{nb}}.
        \]
        Ein Kreis mit gleicher Fläche hätte den Radius
        \[
            \boxed{r=2.36\cdot10^{-4}\,\si{fm}=2.36\cdot10^{-19}\,\si{m}}.
        \]
    \end{answerbox}

    \item \textbf{Plots von $\sigma_{\mathrm{total}}(E_{\mathrm{CM}})$.}

    \begin{formulabox}
        Für die Plots setzen wir $E:=E_{\mathrm{CM}}$ in \si{GeV} ein:
        \[
            \sigma_{\si{nb}}(E)
            =10^7(0.197)^2\frac{(1/137.036)^2}{E^2}\frac{4\pi}{3}
            \bracket{1+\frac{E^4\cdot0.355^2}{\paren{E^2-91.2^2}^2+91.2^2\cdot2.5^2}}.
        \]
        Der Faktor $10^7$ wandelt \si{fm^2} in \si{nb} um.
    \end{formulabox}

    \begin{calculationbox}[title=Lineare Darstellung um die Resonanz]
    \begin{center}
    \begin{tikzpicture}
    \begin{axis}[
        width=0.88\textwidth,
        height=6.0cm,
        xlabel={$E_{\mathrm{CM}}$ in \si{GeV}},
        ylabel={$\sigma_{\mathrm{total}}$ in \si{nb}},
        xmin=50,xmax=150,
        ymin=0,ymax=2.0,
        grid=both,
        samples=500,
        domain=50:150
    ]
    \addplot[thick] {1e7*(0.197)^2*(1/137.036)^2/(x^2)*(4*pi/3)*(1+(x^4*0.355^2)/((x^2-91.2^2)^2+91.2^2*2.5^2))};
    \end{axis}
    \end{tikzpicture}
    \end{center}
    \end{calculationbox}

    \begin{calculationbox}[title=Doppeltlogarithmische Darstellung]
    \begin{center}
    \begin{tikzpicture}
    \begin{axis}[
        width=0.88\textwidth,
        height=6.0cm,
        xlabel={$E_{\mathrm{CM}}$ in \si{GeV}},
        ylabel={$\sigma_{\mathrm{total}}$ in \si{nb}},
        xmode=log,
        ymode=log,
        xmin=10,xmax=500,
        ymin=1e-4,ymax=3,
        grid=both,
        samples=700,
        domain=10:500
    ]
    \addplot[thick] {1e7*(0.197)^2*(1/137.036)^2/(x^2)*(4*pi/3)*(1+(x^4*0.355^2)/((x^2-91.2^2)^2+91.2^2*2.5^2))};
    \end{axis}
    \end{tikzpicture}
    \end{center}
    \end{calculationbox}

    \begin{answerbox}
        Der lineare Plot zeigt einen scharfen Resonanzpeak bei
        \[
            \boxed{E_{\mathrm{CM}}\approx\SI{91.2}{GeV}}.
        \]
        Die doppeltlogarithmische Darstellung zeigt zusätzlich den Verlauf weit unterhalb und oberhalb der Resonanz.
    \end{answerbox}
\end{enumerate}

\subsection{Tafelvortrag zu Aufgabe 1}

\begin{talkbox}
\textbf{Tafelstruktur in 6 Schritten:}
\begin{enumerate}
    \item Start mit der Resonanzbedingung:
    \[
        E_{\mathrm{CM}}=m_Zc^2
        \quad\Rightarrow\quad
        s=m_Z^2c^4.
    \]
    \item Photonterm und $Z$-Term vergleichen:
    \[
        G_\gamma=1,
        \qquad
        G_Z=\frac{s g_eg_\mu}{\Gamma_Z^2}=167.7.
    \]
    \item Winkelabhängigkeit skizzieren:
    \[
        \diffxs\propto1+\cos^2\theta.
    \]
    \item Rutherfordvergleich daneben schreiben:
    \[
        \diffxs_R\propto\frac{1}{\sin^4(\theta/2)}.
    \]
    \item Raumwinkelintegration:
    \[
        \int(1+\cos^2\theta)\dd\Omega=\frac{16\pi}{3}.
    \]
    \item Endergebnisse:
    \[
        \sigma_{\mathrm{total}}=(\hbar c)^2\frac{\alpha^2}{s}G(s)\frac{4\pi}{3},
        \qquad
        \sigma_{\mathrm{peak}}\approx\SI{1.76}{nb},
        \qquad
        r\approx2.36\cdot10^{-4}\,\si{fm}.
    \]
\end{enumerate}
\end{talkbox}

\FloatBarrier
\newpage

%%%%%%%%%%%%%%%%%%%%%%%%%%%%%%%%%%%%%%%%%%%%%%%%%%%%%%%%%%%%%%%%%%%
\section{Aufgabe 2: Solare Neutrinos}

\begin{taskbox}
Die Netto-Reaktion, mit der die Sonne Energie erzeugt, lautet
\[
    4\,{}^1\mathrm{H}\to{}^4\mathrm{He}+2e^+ +2e^-+2\nu_e.
\]
Die entstehenden Positronen annihilieren mit Elektronen. Dadurch wird bei dieser Netto-Reaktion insgesamt
\[
    Q=\SI{26.73}{MeV}
\]
frei. Pro Netto-Reaktion entstehen zwei Elektron-Neutrinos. Die Solarkonstante sei
\[
    S_0=\SI{1361}{W/m^2}.
\]
Weiter gilt für Neutrinos mit typischer Energie $E_\nu=\SI{100}{keV}$ der Wirkungsquerschnitt
\[
    \sigma(\nu_eN)\simeq10^{-21}\,\si{barn}.
\]
\begin{enumerate}[label=(\alph*)]
    \item Berechnen Sie den Fluss der Neutrinos im Abstand der Erde.
    \item Wie viele Neutrinos gehen pro Sekunde durch einen Daumennagel mit Fläche \SI{1}{cm^2}?
    \item Berechnen Sie die Rate von Reaktionen zwischen Neutrinos und Nukleonen bei einer Person mit \SI{70}{kg}, pro Sekunde und pro Jahr.
    \item Berechnen Sie, wie viel Energie durch Elektronen im Körper in einem Jahr frei wird, wenn im Mittel die Hälfte der Neutrinoenergie auf die Elektronen übergeht.
    \item Berechnen Sie daraus die Äquivalentdosis und vergleichen Sie mit der natürlichen jährlichen Strahlenbelastung von ungefähr \SI{1}{mSv/a}.
\end{enumerate}
\end{taskbox}

\subsection{Lernteil: Warum sind solare Neutrinos ungefährlich?}

\begin{learningbox}
Die Sonne produziert enorme Mengen Neutrinos. Deshalb ist der Neutrinofluss an der Erde sehr groß. Trotzdem wechselwirken Neutrinos extrem selten, weil ihr Wirkungsquerschnitt winzig ist:
\[
    \sigma(\nu_eN)\simeq10^{-21}\,\si{barn}=10^{-49}\,\si{m^2}.
\]
Die zentrale Abschätzung ist daher
\[
    R=\Phi_\nu\sigma N_N.
\]
Dabei ist $\Phi_\nu$ der Fluss, $\sigma$ der Wirkungsquerschnitt und $N_N$ die Zahl der Nukleonen im Körper.
\end{learningbox}

\subsection{Abgabe: Ausführliche Lösung}

\begin{enumerate}[label=(\alph*)]
    \item \textbf{Neutrinofluss an der Erde.}

    \begin{calculationbox}
        Die Solarkonstante ist
        \[
            S_0=\SI{1361}{W/m^2}=\SI{1361}{J.s^{-1}.m^{-2}}.
        \]
        Pro Netto-Reaktion wird die Energie $Q=\SI{26.73}{MeV}$ frei und es entstehen zwei Neutrinos. Der Neutrinofluss ist daher
        \[
            \Phi_\nu=2\frac{S_0}{Q}.
        \]
        Zuerst rechnen wir $Q$ in Joule um:
        \[
        \begin{aligned}
            Q
            &=26.73\cdot10^6\,\si{eV}\cdot1.602\cdot10^{-19}\,\si{J/eV}\\
            &=4.283\cdot10^{-12}\,\si{J}.
        \end{aligned}
        \]
        Damit
        \[
        \begin{aligned}
            \Phi_\nu
            &=2\frac{\SI{1361}{J.s^{-1}.m^{-2}}}{4.283\cdot10^{-12}\,\si{J}}\\
            &=6.36\cdot10^{14}\,\si{m^{-2}.s^{-1}}.
        \end{aligned}
        \]
    \end{calculationbox}

    \begin{answerbox}
        Der solare Neutrinofluss an der Erde beträgt
        \[
            \boxed{\Phi_\nu=6.36\cdot10^{14}\,\si{m^{-2}.s^{-1}}}.
        \]
    \end{answerbox}

    \item \textbf{Neutrinos durch einen Daumennagel.}

    \begin{calculationbox}
        Die Fläche sei
        \[
            A=\SI{1}{cm^2}=10^{-4}\,\si{m^2}.
        \]
        Die Zahl der Neutrinos pro Sekunde ist
        \[
            \dot N=\Phi_\nu A.
        \]
        Einsetzen ergibt
        \[
            \dot N
            =6.36\cdot10^{14}\,\si{m^{-2}.s^{-1}}\cdot10^{-4}\,\si{m^2}
            =6.36\cdot10^{10}\,\si{s^{-1}}.
        \]
    \end{calculationbox}

    \begin{answerbox}
        Durch einen Daumennagel mit Fläche \SI{1}{cm^2} gehen pro Sekunde ungefähr
        \[
            \boxed{6.36\cdot10^{10}}
        \]
        solare Neutrinos.
    \end{answerbox}

    \item \textbf{Reaktionsrate im Körper.}

    \begin{calculationbox}
        Der Wirkungsquerschnitt ist
        \[
            \sigma=10^{-21}\,\si{barn}.
        \]
        Mit $\SI{1}{barn}=10^{-28}\,\si{m^2}$ folgt
        \[
            \sigma=10^{-21}\cdot10^{-28}\,\si{m^2}=10^{-49}\,\si{m^2}.
        \]
        Die Zahl der Nukleonen in einer Person mit Masse $m=\SI{70}{kg}$ schätzen wir durch
        \[
            N_N\approx\frac{m}{m_N},
            \qquad
            m_N\approx1.673\cdot10^{-27}\,\si{kg}.
        \]
        Damit
        \[
            N_N\approx\frac{70}{1.673\cdot10^{-27}}=4.19\cdot10^{28}.
        \]
        Die Reaktionsrate ist
        \[
            R=\Phi_\nu\sigma N_N.
        \]
        Einsetzen liefert
        \[
        \begin{aligned}
            R
            &=6.36\cdot10^{14}\,\si{m^{-2}.s^{-1}}
              \cdot10^{-49}\,\si{m^2}
              \cdot4.19\cdot10^{28}\\
            &=2.66\cdot10^{-6}\,\si{s^{-1}}.
        \end{aligned}
        \]
        Pro Jahr ergibt sich mit $\SI{1}{a}=3.156\cdot10^7\,\si{s}$:
        \[
            N_{\mathrm{Jahr}}
            =R\cdot\SI{1}{a}
            =2.66\cdot10^{-6}\cdot3.156\cdot10^7
            =83.9.
        \]
    \end{calculationbox}

    \begin{answerbox}
        Die Rate beträgt
        \[
            \boxed{R=2.66\cdot10^{-6}\,\si{s^{-1}}}.
        \]
        Das entspricht ungefähr
        \[
            \boxed{84\text{ Reaktionen pro Jahr}}.
        \]
    \end{answerbox}

    \item \textbf{Im Körper abgegebene Energie pro Jahr.}

    \begin{calculationbox}
        Die typische Neutrinoenergie ist
        \[
            E_\nu=\SI{100}{keV}.
        \]
        Im Mittel werde die Hälfte auf das Elektron übertragen:
        \[
            E_e=\half E_\nu=\SI{50}{keV}.
        \]
        In Joule ist
        \[
        \begin{aligned}
            E_e
            &=50\cdot10^3\,\si{eV}\cdot1.602\cdot10^{-19}\,\si{J/eV}\\
            &=8.01\cdot10^{-15}\,\si{J}.
        \end{aligned}
        \]
        Aus Teil (c) wurden ungefähr $83.9$ Ereignisse pro Jahr gefunden. Also
        \[
        \begin{aligned}
            \Delta E_{\mathrm{Jahr}}
            &=83.9\cdot8.01\cdot10^{-15}\,\si{J}\\
            &=6.72\cdot10^{-13}\,\si{J}.
        \end{aligned}
        \]
    \end{calculationbox}

    \begin{answerbox}
        Durch die Elektronen wird pro Jahr ungefähr
        \[
            \boxed{\Delta E_{\mathrm{Jahr}}=6.72\cdot10^{-13}\,\si{J}}
        \]
        im Körper abgegeben.
    \end{answerbox}

    \item \textbf{Äquivalentdosis und Vergleich.}

    \begin{calculationbox}
        Für Elektronen gilt der Strahlungswichtungsfaktor
        \[
            w_R=1.
        \]
        Damit ist die Äquivalentdosis
        \[
            H=\frac{\Delta E}{m}w_R.
        \]
        Mit $\Delta E=6.72\cdot10^{-13}\,\si{J}$ und $m=\SI{70}{kg}$ folgt
        \[
        \begin{aligned}
            H
            &=\frac{6.72\cdot10^{-13}\,\si{J}}{\SI{70}{kg}}\\
            &=9.61\cdot10^{-15}\,\si{J/kg}.
        \end{aligned}
        \]
        Da $\SI{1}{Sv}=\SI{1}{J/kg}$ ist,
        \[
            H=9.61\cdot10^{-15}\,\si{Sv/a}.
        \]
        In Millisievert:
        \[
            H=9.61\cdot10^{-12}\,\si{mSv/a}.
        \]
        Verglichen mit einer natürlichen jährlichen Strahlenbelastung von ungefähr $\SI{1}{mSv/a}$ ist
        \[
            \frac{H}{\SI{1}{mSv/a}}=9.61\cdot10^{-12}.
        \]
    \end{calculationbox}

    \begin{answerbox}
        Die jährliche Äquivalentdosis durch solare Neutrinos beträgt
        \[
            \boxed{H=9.61\cdot10^{-15}\,\si{Sv/a}=9.61\cdot10^{-12}\,\si{mSv/a}}.
        \]
        Das ist ungefähr
        \[
            \boxed{10^{11}\text{-mal kleiner}}
        \]
        als eine natürliche jährliche Strahlenbelastung von etwa \SI{1}{mSv/a}.
    \end{answerbox}
\end{enumerate}

\subsection{Tafelvortrag zu Aufgabe 2}

\begin{talkbox}
\textbf{Tafelstruktur in 6 Schritten:}
\begin{enumerate}
    \item Energie pro Netto-Reaktion in Joule umrechnen:
    \[
        Q=\SI{26.73}{MeV}=4.283\cdot10^{-12}\,\si{J}.
    \]
    \item Neutrinofluss aus der Solarkonstante:
    \[
        \Phi_\nu=2\frac{S_0}{Q}=6.36\cdot10^{14}\,\si{m^{-2}.s^{-1}}.
    \]
    \item Durch einen Daumennagel:
    \[
        \dot N=\Phi_\nu\cdot10^{-4}\,\si{m^2}=6.36\cdot10^{10}\,\si{s^{-1}}.
    \]
    \item Körper als Target aus Nukleonen:
    \[
        N_N\approx\frac{70}{1.673\cdot10^{-27}}=4.19\cdot10^{28}.
    \]
    \item Reaktionsrate:
    \[
        R=\Phi_\nu\sigma N_N=2.66\cdot10^{-6}\,\si{s^{-1}}\approx84\,\si{a^{-1}}.
    \]
    \item Dosis:
    \[
        \Delta E\approx84\cdot\SI{50}{keV}=6.72\cdot10^{-13}\,\si{J},
        \qquad
        H=9.61\cdot10^{-15}\,\si{Sv/a}.
    \]
\end{enumerate}
\end{talkbox}

\FloatBarrier
\newpage

%%%%%%%%%%%%%%%%%%%%%%%%%%%%%%%%%%%%%%%%%%%%%%%%%%%%%%%%%%%%%%%%%%%
\section{Gesamtzusammenfassung für die Prüfung}

\begin{summarybox}[title={Formeln, die man können sollte}]
\[
    G(s)=1+\frac{s^2g_eg_\mu}{\paren{s-m_Z^2c^4}^2+m_Z^2c^4\Gamma_Z^2},
    \qquad
    s=E_{\mathrm{CM}}^2.
\]
\[
    \diffxs=(\hbar c)^2\frac{\alpha^2}{s}\frac{G(s)}{4}\paren{1+\cos^2\theta},
    \qquad
    \sigma_{\mathrm{total}}=\int\diffxs\dd\Omega.
\]
\[
    \int(1+\cos^2\theta)\dd\Omega=\frac{16\pi}{3},
    \qquad
    \sigma_{\mathrm{total}}=(\hbar c)^2\frac{\alpha^2}{s}G(s)\frac{4\pi}{3}.
\]
\[
    \Phi_\nu=2\frac{S_0}{Q},
    \qquad
    R=\Phi_\nu\sigma N_N,
    \qquad
    H=\frac{\Delta E}{m}w_R.
\]
\end{summarybox}

\begin{summarybox}[title=Endergebnisse]
\begin{itemize}
    \item $Z$-Beitrag am Resonanzmaximum:
    \[
        \frac{G_Z}{G_\gamma}\approx167.7.
    \]
    \item Winkelabhängigkeit:
    \[
        \diffxs\propto1+\cos^2\theta.
    \]
    \item Totaler Wirkungsquerschnitt am Peak:
    \[
        \sigma_{\mathrm{total}}\approx\SI{1.76}{nb}.
    \]
    \item Kreisradius gleicher Fläche:
    \[
        r\approx2.36\cdot10^{-4}\,\si{fm}=2.36\cdot10^{-19}\,\si{m}.
    \]
    \item Solarer Neutrinofluss an der Erde:
    \[
        \Phi_\nu\approx6.36\cdot10^{14}\,\si{m^{-2}.s^{-1}}.
    \]
    \item Neutrinos pro Sekunde durch \SI{1}{cm^2}:
    \[
        \dot N\approx6.36\cdot10^{10}\,\si{s^{-1}}.
    \]
    \item Reaktionsrate in einer \SI{70}{kg}-Person:
    \[
        R\approx2.66\cdot10^{-6}\,\si{s^{-1}}\approx84\,\si{a^{-1}}.
    \]
    \item Äquivalentdosis:
    \[
        H\approx9.61\cdot10^{-15}\,\si{Sv/a}=9.61\cdot10^{-12}\,\si{mSv/a}.
    \]
\end{itemize}
\end{summarybox}

\begin{warningbox}[title=Typische Fehler]
\begin{itemize}
    \item Beim Resonanzmaximum muss $s=m_Z^2c^4$ gesetzt werden, nicht $s=m_Zc^2$.
    \item In $G(s)$ ist der Photonbeitrag der konstante Term $1$.
    \item Beim totalen Wirkungsquerschnitt darf das Raumwinkelelement $\dd\Omega=\sin\theta\dd\theta\dd\varphi$ nicht vergessen werden.
    \item Für die Umrechnung gilt $\SI{1}{fm^2}=10^7\,\si{nb}$.
    \item Beim Neutrinofluss muss der Faktor $2$ berücksichtigt werden, weil pro Netto-Reaktion zwei Neutrinos entstehen.
    \item Bei der Reaktionsrate muss der Wirkungsquerschnitt von Barn in Quadratmeter umgerechnet werden: $10^{-21}\,\si{barn}=10^{-49}\,\si{m^2}$.
\end{itemize}
\end{warningbox}

\end{document}
